\section{Annotated Verilog Walkthrough}\label{app:verilog}
This appendix explains the final RTL \file{final/rtl/pams\_jqr.v} (SHA-256 \texttt{804fd73a\ldots3760ba0e7f}) region by region.
\begin{itemize}
  \item \textbf{Listings.} Every listing is included verbatim from the file with its original line numbers. Long generated lines are wrapped at the margin (marked $\hookrightarrow$).
  \item \textbf{Generated regions.} Repetitive regions are explained by their generating pattern and one representative instance. The complete file is part of the publication package.
\end{itemize}

\subsection{File organization and generation}\label{app:v-org}
\paragraph{Generation.} The RTL is generated by \file{competitors/models/cgen.py}. It imports the frozen Phase~XI generator \file{models/gen.py}, reusing its circuit builder \sig{Ckt}, the \sig{hpc2} text, the multiplier search \sig{floor\_mult} and the frozen kernels. The builder records:
\begin{itemize}
  \item every masked signal as a list of three share wires plus the pipeline stage at which the value is valid;
  \item every operation as a statement.
\end{itemize}
Whenever two operands of an operation are valid at different stages, the builder inserts share-local delay registers; this is the "aligner". Running the generator with all optimizations disabled reproduces the frozen \file{rtl/pams.v} byte for byte.

\paragraph{Modules.} \Cref{tab:v-modules} lists the modules, and \cref{tab:v-regions} the file regions.
\begin{table}[h]\centering\small
\caption{Modules of \file{pams\_jqr.v} and their dependencies.}\label{tab:v-modules}
\begin{tabular}{@{}llll@{}}
\toprule
module & lines & instances & role \\
\midrule
\sig{hpc2} & 4--19 & 40 in \sig{pams\_jqr\_lane} & 3-share HPC2 AND gadget (\cref{app:v-hpc2}) \\
\sig{pams\_jqr\_lane} & 21--2747 & top & the verification lane \\
\bottomrule
\end{tabular}
\end{table}

\begin{table}[h]\centering\small
\caption{Regions of \file{pams\_jqr.v}.}\label{tab:v-regions}
\begin{tabularx}{\textwidth}{@{}lYl@{}}
\toprule
lines & content & explained in \\
\midrule
4--19 & \sig{hpc2} module & \cref{app:v-hpc2} \\
21--26 & lane ports, stage-1 input registers, mode decode & \cref{app:v-ports} \\
28--45 & public window logic: $\mathrm{qq}$, $r_c$, $r_{c+1}$, $\Low_c$, $\Theta_c$ (stages 1--3) & \cref{app:v-public} \\
46--56 & share-wise modulus switch and anchoring (stages 2--5) & \cref{app:v-front} \\
58--1499 & generated declarations: 199 XOR wires, 11 NOT, 3 AND-with-public, 38 gadget outputs, 1,190 delay registers & \cref{app:v-kernel} \\
1500--1537 & 38 kernel gadget instances & \cref{app:v-gadgets} \\
1538--2729 & 1,190 delay-register assignments (one always block) & \cref{app:v-delays} \\
2730--2746 & control pipeline and ciphertext accumulator & \cref{app:v-acc} \\
\bottomrule
\end{tabularx}
\end{table}

\paragraph{Naming convention of the generated part.} Names are a letter plus a running number:
\begin{itemize}
  \item \sig{x$N$}: share-local XOR;
  \item \sig{n$N$}: NOT, always applied to share 0 only;
  \item \sig{a$N$}: AND of one share with a public bit;
  \item \sig{r$N$}: one-bit delay register;
  \item \sig{g$N$}: 3-bit output of gadget instance \sig{ug$N$}.
\end{itemize}
Masked values have no Verilog vector of their own. Each is three separate wires, one per share; for example, a value with shares (\sig{x120}, \sig{x116}, \sig{x117}) is passed as \sig{\{x117, x116, x120\}} (share 2 first).

\subsection{The HPC2 gadget (lines 4--19)}\label{app:v-hpc2}
\lstinputlisting[firstline=1,lastline=19,firstnumber=1]{\RTL}
\paragraph{Interface.} The gadget computes a 3-share sharing of $c = a \wedge b$ from sharings \sig{a} and \sig{b} with $c_0 \oplus c_1 \oplus c_2 = (a_0 \oplus a_1 \oplus a_2)(b_0 \oplus b_1 \oplus b_2)$. It uses the three fresh bits $r_{01}, r_{02}, r_{12}$ on \sig{r} (line 7). Share $i$ of every signal belongs to domain $i$.

\paragraph{Line by line.}
\begin{itemize}
  \item \emph{Line 10, cycle 1.} \sig{ra}, \sig{rb} and \sig{rr} register $a$, $b$ and $r$.
  \item \emph{Line 11, cycle 1.} \sig{rbr} registers the six refreshed operands $b_j \oplus r_{ij}$ for the ordered pairs $(i, j)$, $i \ne j$, in the order listed on line~8. Each is a function of one share of $b$ and one random bit. Registering it before any use stops glitches from combining $b_j$ with other data.
  \item \emph{Line 12, cycle 2.} The domain-internal products $p_i = a_i b_i$ (registered \sig{p0}).
  \item \emph{Line 13, cycle 2.} \sig{u} $= \bar a_i\, r_{ij}$.
  \item \emph{Line 14, cycle 2.} \sig{v} $= a_i\,(b_j \oplus r_{ij})$.
  \item \emph{Lines 16--18.} Output share $i$ is $c_i = p_i \oplus \bigoplus_{j \ne i} (u_{ij} \oplus v_{ij})$, a combinational XOR of registers of domain $i$ only.
\end{itemize}

\paragraph{Correctness.} For each $i \ne j$, $u_{ij} \oplus v_{ij} = \bar a_i r_{ij} \oplus a_i b_j \oplus a_i r_{ij} = a_i b_j \oplus r_{ij}$. Summing over all $i$, every $r_{ij}$ appears twice (in $c_i$ and $c_j$) and cancels. What remains is $\bigoplus_i a_i b_i \oplus \bigoplus_{i \ne j} a_i b_j = ab$.

\paragraph{Timing.} The latency is two cycles from $a$, $b$ to $c$. A new execution can start every cycle.

\paragraph{Randomness.} The bit $r_{ij}$ is used in cycle 1 (in \sig{rbr}) and again in cycle 2 (through \sig{rr}). Both uses belong to the same execution, so the gadget needs 3 \emph{fresh} bits per execution, i.e.\ per cycle.

\paragraph{Preserved boundary.} The attribute \sig{(* keep\_hierarchy *)} (line 4) stops synthesis from flattening the gadget into surrounding logic, so its internal registers and gate structure survive mapping. \Cref{app:hpc2} shows the mapped circuit.

\subsection{Ports and input registers (lines 21--26)}\label{app:v-ports}
\lstinputlisting[firstline=21,lastline=26,firstnumber=21]{\RTL}
\paragraph{Ports.}
\begin{itemize}
  \item \textbf{Shares (lines 21--23).} \sig{x0}, \sig{x1}, \sig{x2} are the three 12-bit canonical arithmetic shares of the coefficient, $x_i \in [0, 3329)$, which is not checked.
  \item \textbf{Public inputs.} \sig{c} is the received coefficient; 11 bits, of which the low $d$ are used. \sig{mode} selects $d$ (line 24).
  \item \textbf{Frame.} \sig{vin} marks a valid coefficient; \sig{first} and \sig{last} delimit a ciphertext.
  \item \textbf{Randomness.} \sig{rnd} carries the 120 fresh random bits of the cycle.
  \item \textbf{Outputs.} \sig{cdec}/\sig{cvalid} are the per-coefficient decision shares; \sig{dec}/\sig{dvalid} are the ciphertext decision shares.
\end{itemize}

\paragraph{Stage 1 (line 26).} All data inputs are registered, so the lane has no combinational path from an input port. The mode decode $\{0, 1, 2, 3\} \mapsto \{4, 5, 10, 11\}$ produces the 4-bit \sig{d0}.

\subsection{Public window logic (lines 28--45)}\label{app:v-public}
\lstinputlisting[firstline=28,lastline=45,firstnumber=28]{\RTL}
This block computes the public quantities of \cref{eq:low,eq:thetaclosed}; their derivation is in \cref{app:math-anchors,app:math-theta}. Line by line:
\begin{itemize}
  \item \emph{Lines 31--35, \sig{qq}.} Shifts the 12-bit residue left by $e$ into 14 bits and counts how many of $q, 2q, 3q, 4q$ it reaches, i.e.\ $\lfloor r 2^e/q \rfloor$. In gates this is four 14-bit constant comparators and a small adder.
  \item \emph{Line 36.} $e = 13 - d$.
  \item \emph{Line 37.} \sig{rmask} $= 2^{d+1} - 1$.
  \item \emph{Line 38.} \sig{pr} $= (1 - 2c)\,q \bmod 2^{12}$: a constant multiplication in 12-bit arithmetic, whose wrap is the intended mod $2^{12}$.
  \item \emph{Line 39.} $r_c$ = \sig{rc} and $r_{c+1}$ = \sig{rn}, with $2q \bmod 2^{12} = 2562$.
  \item \emph{Line 40.} \sig{qc} and \sig{qn} are $\mathrm{qq}(r_c)$ and $\mathrm{qq}(r_{c+1})$, each 3 bits in $[0, 4]$.
  \item \emph{Lines 42--43, stage 2.} The retiming of variant r registers the intermediate values \sig{qcr}, \sig{qnr}, \sig{e0r}, \sig{c0r} and the shares' \sig{x$i$r}.
  \item \emph{Line 44.} \sig{lowr} $= (2c - 1)2^e + \mathrm{qq}(r_c) \bmod 2^{14}$. The multiplication by $2^e$ is a shift, written as a product.
  \item \emph{Line 45.} \sig{thr} $= 2^{e+1} - 1 + \mathrm{qq}(r_{c+1}) - \mathrm{qq}(r_c)$.
\end{itemize}

\paragraph{Downstream.} \sig{lowr} and \sig{thr} are registered at stage 3 as \sig{lowrq}/\sig{thrq} (line 49), and at stage 4 as \sig{lowp}/\sig{thp} (line 54). \sig{lowp} is consumed at line 56. \sig{thp} becomes \sig{th2} at stage 5, the public threshold bits of the kernel.

\paragraph{Cost.} After synthesis this block (registers plus logic) is $1{,}337\um$, more than in the baseline, whose public logic is a shift and a mask (\cref{tab:resources}). Its slowest stage (\sig{lowrq}, $1.033\ns$ in synthesis STA) is \PAMS{}'s synthesis critical path.

\subsection{Share-wise front end (lines 46--56)}\label{app:v-front}
\lstinputlisting[firstline=46,lastline=56,firstnumber=46]{\RTL}
\paragraph{The modulus switch, \cref{eq:yi}.} The exact switch $\lfloor x_i 2^{14}/q \rfloor = (x_i M) \gg 21$ with $M = 10321339$ was found by the frozen generator's search (\sig{floor\_mult}) as the smallest exact shift.

\paragraph{Stage 3 (line 49).} The 12-bit share is split into its upper and lower 6 bits, and each half is multiplied by $M$ in 48-bit arithmetic: \sig{mh$i$} $= x_{i,11:6} M$ and \sig{ml$i$} $= x_{i,5:0} M$. Since $63 M < 2^{30}$, each product has at most 30 significant bits; synthesis removes the constant-zero upper register bits.

\paragraph{Line 50.} $m_i = 2^6 \cdot$\sig{mh$i$} $+$ \sig{ml$i$} $= x_i M$ exactly.

\paragraph{Line 51.} \sig{s$i$} = bits 34..21 of $m_i$ (14 bits), i.e.\ $y_i$. Bits above 34 are zero because $x_i \le q - 1$ gives $m_i \le (q-1)M < qM = 34{,}359{,}737{,}531 < 2^{35}$.

\paragraph{Stage 4 (lines 53--54).} The switched shares are registered as \sig{s$i$p}. The public constants advance with them.

\paragraph{Stage 5 (lines 55--56).}
\begin{itemize}
  \item \textbf{Anchoring.} \sig{y0} $= \hat y_0 = y_0 - \Low_c + 2 \bmod 2^{14}$: a 14-bit subtract-and-add of public values on share 0 only, \cref{eq:y0}.
  \item \textbf{Other shares.} \sig{y1}, \sig{y2} are registered copies of $y_1$, $y_2$.
  \item \textbf{Threshold.} \sig{th2} = $\Theta_c$.
\end{itemize}

These four 14-bit registers form the interface to the masked kernel.

\paragraph{Locality.} Each \sig{mh$i$}/\sig{ml$i$}/\sig{s$i$p}/\sig{y$i$} depends on share $i$ (and, for $i = 0$, on public data) only. This is what the D-check of \cref{sec:security} verifies on the netlist.

\subsection{Masked kernel: the generating pattern}\label{app:v-kernel}
The kernel evaluates $b = [Z < \Theta_c]$ on the Boolean sharing $(\hat y_0, 0, 0) \oplus (0, y_1, 0) \oplus (0, 0, y_2)$.
\begin{itemize}
  \item \textbf{Source.} It is emitted by \sig{kernel\_pams\_split} in \file{cgen.py}, the frozen kernel \sig{kernel\_pams2} plus just-in-time operand delays.
  \item \textbf{Bit-level functions.} For bit positions $k$:
  \begin{align*}
    S_k &= y_{0,k} \oplus y_{1,k} \oplus y_{2,k} && \text{(no gate: the three shares \emph{are} the three bits)}\\
    T_k &= \mathrm{maj}(y_{0,k-1}, y_{1,k-1}, y_{2,k-1}) && \text{(CSA carry, 1 HPC2)}\\
    z_k &= S_k \oplus T_k \oplus c_k && \text{(sum bit, share-local)}\\
    c_{k+1} &= \mathrm{maj}(S_k, T_k, c_k), \quad c_2 = S_1 \wedge T_1 && \text{(ripple carry, 1 HPC2)}\\
    \lt_{k+1} &= \bigl((z_k \oplus \bar\theta_k) \wedge (\lt_k \oplus \theta_k)\bigr) \oplus \theta_k, \quad \lt_1 = \bar z_0\, \theta_0 && \text{(comparator, 1 HPC2)}
  \end{align*}
  where $\mathrm{maj}(a, b, c) = ((a \oplus b) \wedge (a \oplus c)) \oplus a$ and $\theta_k$ is bit $k$ of $\Theta_c$ (public).
  \item \textbf{Why the comparator is correct.} The recurrence computes $\lt_{k+1} = [Z_{k:0} < \Theta_{c,k:0}]$:
  \begin{itemize}
    \item If $z_k = \theta_k$, the first operand $z_k \oplus \bar\theta_k$ is 1, so $\lt_{k+1} = \lt_k$.
    \item If $z_k \ne \theta_k$, the first operand is 0, so $\lt_{k+1} = \theta_k$, which is 1 iff $z_k = 0 < \theta_k = 1$.
  \end{itemize}
  \item \textbf{Share handling.} All negations and XORs with $\theta_k$ act on share 0 only.
  \item \textbf{Low region.} Bits 1--10 form the ripple, with comparator steps $\lt_2 \ldots \lt_{11}$.
  \item \textbf{Top region.} Bit 11 uses the late carry $c_{11}$. Bits 12--13 use the carry-free zero-test terms $e_k = S_k \oplus T_k \oplus (S_{k-1} \vee T_{k-1})$, with OR $= \overline{\bar x \wedge \bar y}$ realized by one HPC2.
\end{itemize}

\paragraph{Representative instance: bit 0 and bit 1 (lines 58--90).}\leavevmode\par
\lstinputlisting[firstline=58,lastline=78,firstnumber=58]{\RTL}
\begin{itemize}
  \item \textbf{$\lt_1 = \bar S_0 \wedge \theta_0$ (lines 58--61).} With $\theta_0$ public this is linear: share 0 is $\overline{y_{0,0}}\,\theta_0$ (\sig{a102}), shares 1 and 2 are $y_{1,0}\theta_0$ and $y_{2,0}\theta_0$ (\sig{a103}, \sig{a104}).
  \item \textbf{$T_1$ (lines 62--71).} \sig{g111} is the raw output of the majority gadget \sig{ug111} (instance at line 1500). The majority's final "$\oplus\, a$" uses the delayed $y_{0,0}$, registered twice as \sig{r112}$\to$\sig{r113}, giving \sig{x114} $= T_1$ share 0. Shares 1 and 2 of $T_1$ are \sig{g111[1]} and \sig{g111[2]} directly.
  \item \textbf{$z_1 = S_1 \oplus T_1$ (lines 72--74).} \sig{x115}, \sig{x116}, \sig{x117} (shares 0, 1, 2) use the operand bits $y_{i,1}$ delayed by two cycles in \sig{r105}$\to$\sig{r106} (share 0), \sig{r107}$\to$\sig{r108} (share 1) and \sig{r109}$\to$\sig{r110} (share 2).
  \item \textbf{Comparator gadget inputs (lines 77--78).} \sig{x120} $= z_1 \oplus \bar\theta_1$ on share 0 (\sig{r119} is $\bar\theta_1$ delayed twice). \sig{x121} $= \lt_1 \oplus \theta_1$ on share 0. These are the \sig{a}/\sig{b} inputs of the comparator gadget \sig{ug128}.
\end{itemize}

\lstinputlisting[firstline=85,lastline=90,firstnumber=85]{\RTL}
\sig{x133} $= $ \sig{g128[0]} $\oplus\,\theta_1$ (delayed four times in \sig{r129}..\sig{r132}) is share 0 of $\lt_2$. The same three-gadget pattern (majority $T_k$, comparator $\lt_{k+1}$, carry $c_{k+1}$) repeats for $k = 2, \dots, 10$, each time with longer delay chains.

\subsection{Gadget instances (lines 1500--1537)}\label{app:v-gadgets}
\lstinputlisting[firstline=1500,lastline=1505,firstnumber=1500]{\RTL}
\lstinputlisting[firstline=1530,lastline=1537,firstnumber=1530]{\RTL}

\paragraph{Instance order.} The generator creates gadgets in the order the kernel needs them, and gadget $g$ receives \sig{rnd[$3g+2$:$3g$]}. Read from the file, the instances fall into three groups, so every random bit is wired to exactly one gadget:
\begin{itemize}
  \item \emph{The ripple region.} Instances \sig{ug111}, \sig{ug128}, \sig{ug134} (bit 1), \sig{ug153}, \sig{ug173}, \sig{ug187} (bit 2), \ldots, \sig{ug1249}, \sig{ug1301}, \sig{ug1331} (bit 10) form ten (majority, comparator, carry) triples, using \sig{rnd[89:0]}.
  \item \emph{The top region.}
    \begin{itemize}
      \item \sig{ug1341}, \sig{ug1417}, \sig{ug1449}: the majorities $T_{11}$, $T_{12}$, $T_{13}$;
      \item \sig{ug1438}, \sig{ug1470}: the ORs $S_{11} \vee T_{11}$ and $S_{12} \vee T_{12}$;
      \item \sig{ug1483}: the OR tree of $e_{12}$, $e_{13}$;
      \item \sig{ug1535}: the AND $\bar e_{11} \wedge \overline{(e_{12} \vee e_{13})}$, i.e.\ $[Z_{13:11} = 0]$;
      \item \sig{ug1542}: the final AND with $\lt_{11}$, producing $b$ = \sig{cdec}.
    \end{itemize}
    These use \sig{rnd[113:90]}.
  \item \emph{The accumulator.} \sig{uacc} and \sig{ufin} use \sig{rnd[119:114]}.
\end{itemize}

\paragraph{Constant shares.} Majority gadgets such as \sig{ug111} receive the literal \sig{1'b0} in one share of each operand (operands $(y_0, y_1, 0)$ and $(y_0, 0, y_2)$). In the routed netlist these become tie cells inside the gadget hierarchy. They are part of the verified structure: SILVER and PROLEAD tested the gadget with all share inputs free, so constant shares are a special case.

\subsection{Delay registers (lines 1538--2729)}\label{app:v-delays}
\lstinputlisting[firstline=1538,lastline=1560,firstnumber=1538]{\RTL}
\paragraph{Pattern.} Each line is a one-bit register \sig{r$N$ <= }\emph{source}. Chains of such registers delay one share bit, a public $\theta_k$ bit, or one share of a gadget output until its consumer's stage. For example, line 1555 onward holds $\theta_1$ for four cycles.

\paragraph{How the chains arise.} The aligner creates a new chain per consumer. The RTL therefore declares 1,190 delay registers. Many chains are duplicates: the same source delayed by the same depth for two consumers.

\paragraph{Counts after synthesis.} Resolving every register to its (source, depth) pair gives exactly 782 distinct registers. That is precisely the number of alignment registers in the synthesized netlist: Yosys merges registers with identical inputs. Of these:
\begin{itemize}
  \item 406 delay single-domain operand bits $y_{i,k}$, the just-in-time CSA operands;
  \item 240 delay public $\theta_k$ bits;
  \item 100 delay share-local intermediate wires;
  \item 36 delay one share of a gadget output.
\end{itemize}
They cost $3{,}558\um$ including the little logic in front of them (\cref{tab:resources}). They are the second-largest block after the gadgets.

\paragraph{Locality.} A one-bit register holding share $i$ of something is domain $i$. One holding a $\theta$ bit is public.

\subsection{Control pipeline and accumulator (lines 2730--2746)}\label{app:v-acc}
\lstinputlisting[firstline=2730,lastline=2746,firstnumber=2730]{\RTL}
\paragraph{Line by line.}
\begin{itemize}
  \item \emph{Lines 2731--2732.} 32-bit shift registers delay \sig{vin}, \sig{first} and \sig{last} (\sig{f2p} is an unused duplicate that synthesis removes). Tap 30 corresponds to stage 31, because the input register adds one stage.
  \item \emph{Line 2733.} \sig{cdec} is the output sharing of \sig{ug1542}; \sig{cvalid} marks it.
  \item \emph{Lines 2734--2735.} \sig{init\_p} is high for the first two coefficients of a ciphertext.
  \item \emph{Lines 2738--2740.} \sig{acc\_in} is the trivial sharing $(1, 0, 0)$ of 1 when a chain starts, and otherwise the gadget's own output \sig{acc\_out}. \sig{uacc} computes \sig{acc\_out} = \sig{acc\_in} $\wedge\, b$ two cycles later. Because of the two-cycle latency, even- and odd-indexed coefficients form two independent product chains.
  \item \emph{Lines 2741--2743.} \sig{ufin} ANDs the two chains: \sig{acc\_last} is the previous chain's value.
  \item \emph{Lines 2745--2746.} \sig{dvalid} rises four cycles after the last coefficient's \sig{cvalid}, when \sig{dec} holds $D = \bigwedge b$.
\end{itemize}

\paragraph{Timing behaviour.} \sig{cvalid} and \sig{dvalid} are functions of \sig{vin}/\sig{first}/\sig{last} only. Phase XIV verified on all final netlists that they depend on no share and no random bit.

\subsection{Reference tables}\label{app:v-tables}
\begin{table}[h]\centering\small
\caption{Glossary of the important signals of \sig{pams\_jqr\_lane}.}\label{tab:v-glossary}
\begin{tabularx}{\textwidth}{@{}llrlX@{}}
\toprule
signal & lines & width & kind & meaning \\
\midrule
\sig{x$i$}, \sig{x$i$q}, \sig{x$i$r} & 21, 25--26, 42--43 & 12 & share $i$ & arithmetic share $x_i$ (stages 0, 1, 2) \\
\sig{c0}, \sig{d0} & 25--26 & 11, 4 & public & registered $c$; decoded $d$ \\
\sig{e0}, \sig{rmask}, \sig{pr}, \sig{rc}, \sig{rn} & 36--39 & 4, 12 & public & $e = 13-d$; $2^{d+1}-1$; $(1-2c)q \bmod 2^{12}$; $r_c$; $r_{c+1}$ \\
\sig{qc}, \sig{qn} $\to$ \sig{qcr}, \sig{qnr} & 40, 42--43 & 3 & public & $\mathrm{qq}(r_c)$, $\mathrm{qq}(r_{c+1})$ \\
\sig{lowr}$\to$\sig{lowrq}$\to$\sig{lowp} & 44, 48--49, 53--54 & 14 & public & $\Low_c$ (stages 2c, 3, 4) \\
\sig{thr}$\to$\sig{thrq}$\to$\sig{thp}$\to$\sig{th2} & 45, 48--49, 53--56 & 14 & public & $\Theta_c$ (stage 5 = \sig{th2}) \\
\sig{mh$i$}, \sig{ml$i$} & 48--49 & 48 & share $i$ & partial products $x_{i,11:6} M$, $x_{i,5:0} M$ (stage 3) \\
\sig{m$i$}, \sig{s$i$} & 50--51 & 48, 14 & share $i$ & $x_i M$; $y_i = F_{14}(x_i)$ (combinational) \\
\sig{s$i$p} & 53--54 & 14 & share $i$ & $y_i$ (stage 4) \\
\sig{y0}, \sig{y1}, \sig{y2} & 55--56 & 14 & share 0/1/2 & $\hat y_0$, $y_1$, $y_2$: kernel operands (stage 5) \\
\sig{x$N$}, \sig{n$N$}, \sig{a$N$} & 58--1499 & 1 & one share & share-local XOR / NOT / AND-with-public \\
\sig{r$N$} & 58--1499, 1538--2729 & 1 & one share or public & delay register \\
\sig{g$N$} & 58--1499 & 3 & 3 shares & output of gadget \sig{ug$N$} \\
\sig{rnd} & 22 & 120 & random & fresh bits; \sig{rnd[$3g+2$:$3g$]} to gadget $g$ \\
\sig{cdec}, \sig{cvalid} & 2733 & 3, 1 & 3 shares, public & $b$ (stage 31) \\
\sig{acc\_in}, \sig{acc\_out}, \sig{acc\_last} & 2739--2742 & 3 & 3 shares & accumulator chains \\
\sig{dec}, \sig{dvalid} & 2743--2746 & 3, 1 & 3 shares, public & $D$ (4 cycles after the last \sig{cvalid}) \\
\bottomrule
\end{tabularx}
\end{table}

\begin{table}[h]\centering\small
\caption{Register and pipeline-stage map (stage $s$ = clock edges after the coefficient is applied).}\label{tab:v-stages}
\resizebox{\textwidth}{!}{%
\begin{tabular}{@{}rlll@{}}
\toprule
stage & share lanes ($\times 3$) & public lane & kernel / control \\
\midrule
1 & \sig{x$i$q} (12 b) & \sig{c0} (11 b), \sig{d0} (4 b) & \sig{vp[0]}, \sig{fp[0]}, \sig{lp[0]} \\
2 & \sig{x$i$r} (12 b) & \sig{qcr}, \sig{qnr} (3 b), \sig{e0r} (4 b), \sig{c0r} (11 b) & \\
3 & \sig{mh$i$}, \sig{ml$i$} ($\le$ 30 significant bits each) & \sig{lowrq}, \sig{thrq} (14 b) & \\
4 & \sig{s$i$p} (14 b) & \sig{lowp}, \sig{thp} (14 b) & \\
5 & \sig{y$i$} (14 b) & \sig{th2} (14 b) & kernel inputs; first gadget layer \\
5--31 & delay chains of $y_{i,k}$ & delay chains of $\theta_k$ & 13 gadget levels of 2 cycles \\
31 & & & \sig{cdec}, \sig{cvalid} \\
33--35 & & & \sig{uacc}, \sig{ufin}; \sig{dec}, \sig{dvalid} at stage 35 after the last coefficient \\
\bottomrule
\end{tabular}}
\end{table}

\begin{table}[h]\centering\small
\caption{Source lines, equations and figures.}\label{tab:v-lines}
\begin{tabularx}{\textwidth}{@{}lXll@{}}
\toprule
lines & computes & equation & figure \\
\midrule
21--26 & input shares, public $c$, $d$ & \cref{eq:target} & \cref{fig:datapath} (left) \\
31--35, 40 & $\mathrm{qq}(r) = \lfloor r 2^e/q \rfloor$ & \cref{eq:app-low} & \cref{fig:datapath} (public lane) \\
36--39 & $e$, $r_c$, $r_{c+1}$ & \cref{eq:low,eq:thetaclosed} & \cref{fig:datapath} \\
44 & $\Low_c$ & \cref{eq:low} & \cref{fig:datapath} \\
45 & $\Theta_c$ & \cref{eq:thetaclosed}, \cref{eq:theta} & \cref{fig:datapath} \\
49--51, 54 & $y_i = F_{14}(x_i)$ & \cref{eq:yi} & \cref{fig:datapath}; \cref{fig:placement} (orange) \\
56 & $\hat y_0 = y_0 + 2 - \Low_c$ & \cref{eq:y0} & \cref{fig:datapath} \\
58--1537 & $Z$ and $[Z < \Theta_c]$: CSA, ripple, comparator, zero test & \cref{eq:Z,eq:b} & \cref{fig:slice}; \cref{fig:datapath} (kernel) \\
4--19 & HPC2 AND & -- & \cref{fig:hpc2,fig:hpc2-mapped} \\
1538--2729 & operand alignment & -- & \cref{fig:placement} (green) \\
2730--2746 & $D = \bigwedge b$, valid flags & \cref{eq:pred} & \cref{fig:datapath} (right) \\
\bottomrule
\end{tabularx}
\end{table}

\subsection{Tracing one coefficient through the RTL}\label{app:v-trace}
\Cref{tab:trace-e1} follows example E1 through the pipeline of the final RTL, one register stage at a time.

\paragraph{How the trace was made.} \file{paper/scripts/appendix\_data.py} simulates \file{final/rtl/pams\_jqr.v} with Icarus Verilog and applies E1 alone:
\begin{itemize}
  \item E1 is presented at one falling clock edge with \sig{vin}, \sig{first} and \sig{last} all high, so the ciphertext has a single coefficient.
  \item Every other cycle carries a filler coefficient ($x = 0$, $c = 1$, $d = 4$) whose decision is 0. A decision of 1 at stage 31 can therefore only belong to E1.
  \item \sig{rnd} is fresh every cycle, as the security argument requires.
\end{itemize}
The script reads the registers at each stage, after the given number of rising edges. It asserts each value against the reference in \file{data/examples.json}, which is computed independently from \file{models/pams.py}:
\begin{itemize}
  \item the registered shares;
  \item the six partial products $\lfloor x_i/64 \rfloor M$ and $(x_i \bmod 64) M$;
  \item $F_{14}(x_i)$ and the aligned operands;
  \item $\mathrm{qq}(r_c)$, $\mathrm{qq}(r_{c+1})$, $e$, $\Low_c$ and $\Theta_c$;
  \item the first CSA majority;
  \item the decision.
\end{itemize}
It also asserts that \sig{cvalid} is high at stage 31 and at no other stage.

\begin{table}[h]\centering\small
\caption{Example E1 ($x = 3300$, shares $(487, 1308, 1505)$, $d = 4$, $c = 0$) in the registers of \file{pams\_jqr.v}, by pipeline stage. "Recombined" values are the XOR of a gadget's three output shares, formed only in simulation for this table; the hardware never computes them.}\label{tab:trace-e1}
\begin{tabularx}{\textwidth}{@{}rlR{6.2cm}Y@{}}
\toprule
stage & register(s) & value & meaning \\
\midrule
1 & \sig{x$i$q}; \sig{c0}, \sig{d0} & \texttt{487, 1308, 1505; 0, 4} & registered shares $x_i$; $c$, $d$ \\
2 & \sig{x$i$r}; \sig{qcr}, \sig{qnr}, \sig{e0r} & \texttt{487, 1308, 1505; 0, 4, 9} & $x_i$; $\mathrm{qq}(r_c)$, $\mathrm{qq}(r_{c+1})$, $e$ \\
3 & \sig{mh$i$}, \sig{ml$i$} & \texttt{72249373, 402532221, 206426780, 288997492, 237390797, 340604187} & $\lfloor x_i/64\rfloor M$, $(x_i \bmod 64) M$ for $i = 0, 1, 2$ \\
3 & \sig{lowrq} & \texttt{15872} & $\Low_c$ (Equation~(7)) \\
4 & \sig{s$i$p}; \sig{lowp} & \texttt{2396, 6437, 7407; 15872} & $F_{14}(x_i)$; $\Low_c$ \\
5 & \sig{y$i$}; \sig{th2} & \texttt{2910, 6437, 7407; 1027} & $\hat y_0, y_1, y_2$; $\Theta_c$ \\
7 & \sig{g111} (recombined) & \texttt{1} & $T_1 = \mathrm{maj}$ of bit 0 of the three operands \\
31 & \sig{g1542} (recombined); \sig{cvalid} & \texttt{1; 1} & $b = [Z < \Theta_c]$; decision marker \\
31 & \sig{cdec} shares $(b_0, b_1, b_2)$ & \texttt{0, 1, 0} & $b_0 \oplus b_1 \oplus b_2 = 1$ \\
\bottomrule

\end{tabularx}
\end{table}

\paragraph{Reading the trace.}
\begin{itemize}
  \item \emph{Stages 1--2.} The shares pass through two register stages unchanged. In parallel the public lane derives $e = 13 - d = 9$, $\mathrm{qq}(r_0) = 0$ and $\mathrm{qq}(r_1) = 4$ from $(c, d) = (0, 4)$; here $r_0 = 3329 \bmod 32 = 1$ and $r_1 = 31$.
  \item \emph{Stage 3.} Each share is split as $x_i = 64 \lfloor x_i/64 \rfloor + (x_i \bmod 64)$ and both halves are multiplied by $M = 10321339$. For share 0: $487 = 64 \cdot 7 + 39$, so the products are $7M = 72249373$ and $39M = 402532221$. The public anchor is $\Low_0 = (2 \cdot 0 - 1) 2^9 + 0 \equiv 15872$.
  \item \emph{Stage 4.} Recombining, $m_0 = 64 \cdot 72249373 + 402532221 = 5026492093$, and $m_0 \gg 21 = 2396 = F_{14}(487)$.
  \item \emph{Stage 5.} Share 0 is anchored: $2396 - 15872 + 2 \equiv 2910$. The threshold $\Theta_0 = 2^{10} - 1 + 4 - 0 = 1027$ is registered in \sig{th2}.
  \item \emph{Stage 7.} The first gadget layer has finished: the majority of bit 0 of $(2910, 6437, 7407)$, i.e.\ of $(0, 1, 1)$, is 1.
  \item \emph{Stage 31.} The final AND gadget produces $b = 1$ ($Z = 370 < 1027$). Its three shares $(0, 1, 0)$ are randomized by the gadget's fresh bits; only their XOR is the decision.
\end{itemize}
The stages between 7 and 31 hold the remaining gadget layers and the operand alignment of \cref{app:v-delays}. They are not shown because no single register there holds a named mathematical quantity in the clear.

\subsection{The baseline RTL for comparison}\label{app:v-base}
\file{final/rtl/base\_jq.v} (3,167 lines) has the same structure, ports and accumulator. Its 171-bit \sig{rnd} feeds 57 gadgets. Its front end differs:
\lstinputlisting[firstline=28,lastline=46,firstnumber=28]{\RTLB}
\paragraph{Differences from \PAMS{}.}
\begin{itemize}
  \item \emph{Public offset.} The public offset $(2c - 1)2^{12} - 1$ (line 31) is a shift: there is no division-like window logic.
  \item \emph{Wider shares.} The share-wise constant is the 34-bit $M = 10569051393$ (line 38), giving a 24-bit $\lfloor x_i 2^{24}/q \rfloor$ (line 40).
  \item \emph{Mode shift.} A mode-dependent right shift by $11 - d$ (line 44) aligns the $d$ significant bits. The 24-bit operands then enter a 24-bit kernel: 23 CSA majorities, 12 ripple carries to bit 13, and a carry-free zero test of bits 13--23 enabled by \sig{en2}.
\end{itemize}
